% Requires \usepackage{booktabs,longtable}
\begin{longtable}{p{0.24\textwidth} p{0.13\textwidth} p{0.53\textwidth}}
\caption{Mutational-context covariates used in the context-control analysis.}\label{tab:context_covariates}\\
\toprule
Covariate & Encoding & Definition \\
\midrule
\endfirsthead
\toprule
Covariate & Encoding & Definition \\
\midrule
\endhead
Reference codon & Categorical & Wild-type DNA codon (U converted to T). \\
Mutant codon & Categorical & Mutant DNA codon (U converted to T). \\
Mutated codon position & Categorical & Position 1, 2, or 3 within the codon. \\
Reference nucleotide & Categorical & Wild-type nucleotide at the mutated position. \\
Mutant nucleotide & Categorical & Mutant nucleotide at the mutated position. \\
Nucleotide substitution & Categorical & Directed substitution, for example C>T. \\
Transition status & Numeric & 1 for A>G, G>A, C>T, or T>C; otherwise 0. \\
Transversion status & Numeric & 1 minus transition status. \\
Reference codon GC content & Numeric & Fraction of G or C nucleotides in the reference codon. \\
Mutant codon GC content & Numeric & Fraction of G or C nucleotides in the mutant codon. \\
Change in codon GC content & Numeric & Mutant minus reference codon GC content. \\
Reference within-codon CpG status & Numeric & 1 if the reference codon contains CG; otherwise 0. \\
Mutant within-codon CpG status & Numeric & 1 if the mutant codon contains CG; otherwise 0. \\
Change in within-codon CpG status & Numeric & Mutant minus reference within-codon CpG status. \\
Local codon degeneracy & Categorical & Number of possible nucleotides at the mutated position that preserve the reference amino acid. \\
\bottomrule
\end{longtable}
